\begin{abstract}
Deep learning has transformed medical image analysis, becoming central to many clinical applications. However, real-world deployment remains challenging due to computational demands, latency constraints, and privacy concerns associated with cloud-based processing. Despite these limitations, state-of-the-art models increasingly favor larger architectures and datasets, reinforcing a scale-driven paradigm that is often impractical for resource-constrained systems.
This study systematically evaluates deep learning models with fewer than one million parameters on brain MRI classification and histological glioma grading, focusing on edge efficiency and cross-dataset generalization. On saturated classification benchmarks, compact models achieve diagnostic accuracy statistically comparable to architectures an order of magnitude larger; the top-performing compact model reaches 98.29\% accuracy at 0.36 GFLOPs, executing six times faster than ResNet-50 with less than a one-percentage-point difference in accuracy. 
However, a performance gap remains on more challenging cross-dataset tasks, with errors concentrated in low-confidence predictions sensitive to input perturbations. The most efficient architecture, which uses parallel atrous convolution branches to capture multi-scale features, is particularly vulnerable to scale-sensitive inputs.
To address this limitation, we propose \textbf{MCG-TTA}, an architecture-informed inference strategy designed to improve robustness while limiting computational overhead. MCG-TTA improves slice-level F1 by an average of 1.43 percentage points (and up to 1.83 percentage points) at approximately 1.9$\times$ the baseline latency, while retaining a model size of 0.91 million parameters. These findings suggest that, for resource-constrained clinical applications, uncertainty-aware inference can recover residual accuracy more efficiently than increasing model complexity.
\end{abstract}

\begin{IEEEkeywords}
Magnetic resonance imaging (MRI), brain tumor classification, compact neural networks, edge computing, test-time augmentation, glioma grading, uncertainty estimation.
\end{IEEEkeywords}

